\documentclass[10pt,a4paper]{amsart}
\usepackage[margin=19mm]{geometry}
\usepackage{amsmath,amssymb,mathtools}
\usepackage[scr=rsfs,cal=euler]{mathalfa}
\usepackage{xfrac}
\usepackage[unicode,hidelinks,bookmarks=false]{hyperref}
\newcommand{\ie}{i.e.\ }
\newcommand{\cf}{cf.\ }
\newcommand{\resp}{respectively\ }
\newcommand{\Subsection}{\subsection}
\providecommand{\nobreakdash}{\nobreak\hbox{-}}
\setlength{\emergencystretch}{2em}
\multlinegap0pt
\usepackage{dsfont}
\usepackage{amssymb}
\usepackage{mathtools}
\usepackage{float}
\usepackage{xcolor}
\newtheorem{defi}{Definition}
\newtheorem{theorem}{Theorem}
\newcommand{\PP}{{\rm I\! P}}
\newcommand{\cX}{\mathcal{X}}
\newcommand{\cZ}{\mathcal{Z}}
\newcommand{\mf}[1]{\mathfrak{#1}}
\title{Universal computation with magic Hamiltonians}
\author{Marius Junge, Jason Pollack, Luke Visser}
\date{Source: arXiv:2609.14757v1 (CC BY 4.0)}
\begin{document}
\maketitle

\noindent\emph{Source and licence.} Universal computation with magic Hamiltonians. Version arXiv:2609.14757v1 (CC BY 4.0), distributed by arXiv under the Creative Commons Attribution 4.0 International licence, http://creativecommons.org/licenses/by/4.0/, as recorded in the arXiv licence metadata for this exact version and shown on the abstract page https://arxiv.org/abs/2609.14757v1. Full licence notice: \texttt{licenses/arxiv-2609.14757-CC-BY-4.0.txt}.\par\smallskip

\noindent\textit{Editorial scope.} Counted unit: the article's theorem thm:commutators\_algorithmically (source0.tex lines 964--970). The article's complete original proof of it is delivered with it (source0.tex lines 972--1027, one \texttt{proof} environment of 56 lines). No mathematics is written, repaired or re-derived: the statement and the proof are the article's own lines, in the article's own notation, and no retained line is substituted or re-flowed.  Retained uncounted as the source-local context the proof needs: lines 447--456; lines 688--702; lines 1408--1408.  Nothing was omitted from any retained span, so the package carries no omission record.  No retained line contains a citation, so the wrapper carries no bibliography and no bibliography entry.  No retained line contains a figure environment, an image-inclusion command or drawing code, so no image is delivered and the package declares no asset dependency.  Editorial formatting only, in addition: an AMS \texttt{amsart} preamble that restates the article's own declarations and macros used by the retained lines, namely the environments \texttt{defi}, \texttt{theorem} on separate counters, together with the standard packages those lines need.  The retained environments print in this document as Definition 1, Theorem 1 in order of appearance, whereas the article prints the same items with the numbering of its own full document.  The complete native source of the article as distributed in the arXiv e-print for this exact version is bundled unchanged as \texttt{sources/210353/source0.tex}.
\begin{defi}\label{def:length_function}
    Define the length function $l(R)$ as the minimal number of commutators with elements from a generator set $G$ that is required to make $R$. If we define $G_k$ iteratively, with $G_1 = G$ and $G_{k+1}$ as 
    \begin{equation}
        G_{k+1} = G_k \cup \{[P_i, P_j] ~:~ P_i \in G_1,~ P_j \in G_k \},
    \end{equation}
    we can formally define $l(R)$ as
    \begin{equation}
        l(R)=\inf \big\{ k ~:~ R \in \operatorname{span}(G_k)\}.
    \end{equation}
\end{defi}

\begin{theorem}\label{thm:XdoubleZsingle_extensions}
    The Lie algebra generated by $\cX(2)$ and $\cZ(1)$ combined with $X_1$ and/or $Z_1 Z_2$ is given by
    \begin{equation*}
    \begin{aligned}
        i) && \langle \cX(2) \cup \cZ(1) \cup \{X_1\} \rangle =& \{Z_j ~:~ j \in [n]\} \cup \{X_1, Y_1\} \\ 
        &&& \cup \{\alpha_i Z^{\otimes m} \beta_{i+m+2} ~:~ \alpha, \beta \in \{X,Y\},~ m \in [n-2] \} \\
        &&& \cup \{Z^{\otimes m} \beta_{m+1} ~:~ \beta \in \{X,Y\}, m \in [n-1] \}  \\
        &&& = \mf{so}(2n+1)\\
        ii) && \langle \cX(2) \cup \cZ(1) \cup \{Z_1 Z_2\} \rangle =& \langle \cX(2) \cup \cZ(2) \cup \cZ(1) \rangle  \\
        && =& S_{z_n} \backslash \{(Z_1 Z_2 \cdots Z_n), I\} = \oplus_{i=1}^2 \mf{su}(2^{n-1}), \\
        iii) && \langle \cX(2) \cup \cZ(1) \cup \{X_1, Z_1 Z_2\} \rangle =& \PP_n \backslash \{I \} = \mf{su}(2^n)
    \end{aligned}
    \end{equation*}
    with $S_{z_n}$ all Pauli strings with an even number of Pauli $X$ and Pauli $Y$ operators, and $\PP_n$ all possible Pauli strings. The dimensions of the generated Lie algebras are respectively $2n^2+n$, $2( 4^{n-1}-1)$, and $4^n-1$ for $n\geq 3$.
\end{theorem}

\section{Pseudo-algorithm}\label{app:alg-computation}

\begin{theorem}\label{thm:commutators_algorithmically}
    In the settings of Theorem \ref{thm:XdoubleZsingle_extensions}, any Pauli string can be written as a chain of commutators using the generators
    \begin{equation}
        \langle \cX(2) \cup \cZ(1) \cup \{X_1\} \cup \{Z_1 Z_2\} \rangle = P_n \backslash \{I \},
    \end{equation}
    where the chain has length at most $4n^3$ and can be found algorithmically.
\end{theorem}

\begin{proof}
We provide an explicit algorithm, which we build using three local subroutines. Each subroutine starts with a specific type of Pauli string and applies a sequence of commutators with elements of the generator set to change the Pauli string. We will denote $W_{\leq k}$ for an arbitrary Pauli string on indices $1,2,\ldots, k$ and $W_{\geq k}$ for an arbitrary Pauli string on indices $k, k+1, \ldots, n$. The identity operator $I$ acts on any indices not mentioned.

The three routines are
\begin{equation*}
\begin{aligned}
    1)&& \quad W_{\leq k-1} Z_k I_{k+1} W_{\geq k+2} &\to W_{\leq k-1} I_k Z_{k+1} W_{\geq k+2}, \\
    2)&& \quad Z_1 I_2 W_{\geq 3} & \to Z_1 Z_2 W_{\geq 3}, \\
    3) &&\quad Z_1 W_{\geq k+1} & \to Z_1 X_k W_{\geq k+1}.
\end{aligned}
\end{equation*}
The first subroutine interchanges a $Z$ operator in a Pauli chain with a neighbouring $I$ operator and the second subroutine adds an extra $Z$ operator at the beginning of the Pauli string. These two subroutines together allow us to create any Pauli string consisting of only $I$ and $Z$. The third subroutine makes a $X$ or $Y$ at position $k$ without touching any of the elements past index $k$.

These three subroutines together allow us to make any Pauli string by starting with $Z_1$ and then creating any Pauli string iteratively from highest to lowest index. If the Pauli string has leading $I$'s, we can get rid of the final $Z_1$ as well by either simply shifting it over with the first subroutine if the leading non-identity $\alpha_k = Z$ or using a slightly modified version of the third subroutine that still starts with $Z_1$ but makes $X_k$ instead of $Z_1 X_k$. If there are no leading $I$'s, the string on the first two indices $\alpha_1 \alpha_2$ can be fixed using commutators with $X_1, Z_1, Z_2, X_1 X_2,$ and $Z_1 Z_2$, as these operators fully control the first two qubits.

The first subroutine starts with $W_{\leq k-1} Z_k I_{k+1} W_{\geq k+2}$. Taking the commutators
\begin{equation}
    Z_k I_{k+1} \stackrel{X_k X_{k+1}}{\longleftrightarrow} Y_k X_{k+1} \stackrel{Z_k}{\longleftrightarrow} X_k X_{k+1} \stackrel{Z_{k+1}}{\longleftrightarrow}  X_k Y_{k+1} \stackrel{X_k X_{k+1}}{\longleftrightarrow} I_k Z_{k+1},
\end{equation}
the subroutine shifts one $Z$ one index at the cost of 4 commutations, which is the same as adding 4 to the length function $l$ defined in Definition \ref{def:length_function}.

The second subroutine starts with $Z_1 I_2 W_{\geq 3}$ and uses a sequence of commutations
\begin{equation}
    Z_1 I_2 W_{\geq 3} \stackrel{X_1}{\longleftrightarrow} Y_1 I_2 W_{\ge 3} \stackrel{Z_1 Z_2}{\longleftrightarrow} X_1 Z_2 W_{\ge 3} \stackrel{Z_1}{\longleftrightarrow}Y_1 Z_2 W_{\ge 3} \stackrel{X_1}{\longleftrightarrow} Z_1 Z_2 W_{\geq 3}.
\end{equation}
to increase the number of $Z$ matrices in the Pauli chain by 1 at the cost of 4 commutations. Together with the first subroutine, this allows us to make any string consisting of only $I$ and $Z$. The cost of creating $Z_k$ instead of $I_k$ is 4 commutations to make an extra $Z$ at index 2, and $4(k-2)$ to shift the index over from $Z_2$ to $Z_k$, for a total of $4(k-1)$ commutations.

The third and final subroutine creates an element $Z_1 Y_k W_{\geq k+1}$ or $Z_1 X_k W_{\geq k+1}$ from $Z_1 W_{\geq k+1}$. 
The subroutine starts with $Z_1 I_2 Z_3 Z_4 \cdots Z_k W_{\geq k+1}$, which can be made from $Z_1 W_{\geq k+1}$ using the previous two subroutines for $\sum_{j=3}^k 4(k-1) = 2k(k-1)-4$ commutations. Then the subroutine adds $X$ to the second index via the commutators
\begin{equation}
    Z_1 I_2 Z_3 Z_4 \cdots Z_k W_{\geq k+1} \stackrel{X_1 X_2}{\longleftrightarrow} Y_1 X_2 Z_3 Z_4 \cdots Z_k W_{\geq k+1} \stackrel{X_1}{\longleftrightarrow} Z_1 X_2 Z_3 Z_4 \cdots Z_k W_{\geq k+1}.
\end{equation}
After that, the subroutine shifts the $X$ to the end of the chain of $Z$'s by eating up the $Z$ matrices. This is implemented by commuting with $X_j X_{j+1}$ and $Z_{j+1}$ for $j = 2, 3 \cdots k-1$ to get
\begin{equation}
\begin{aligned}
        Z_1 X_2 Z_3 \cdots Z_kW_{\geq k+1}  &\stackrel{X_2 X_3}{\longleftrightarrow} Z_1 I_2 Y_3 \cdots Z_kW_{\geq k+1}  \stackrel{Z_3}{\longleftrightarrow} Z_1 I_2 X_3 Z_4 \cdots Z_kW_{\geq k+1} \\
        &\cdots  \\
        &\stackrel{X_{k-1} X_k}{\longleftrightarrow} Z_1 I_2 \cdots I_{k-1} Y_k W_{\geq k+1} \stackrel{Z_k}{\longleftrightarrow} Z_1 I_2 \cdots I_{k-1} X_k W_{\geq k+1}. 
\end{aligned}
\end{equation}
Note that applying the last $Z_k$ instead yields  $Z_1 Y_k W_{\geq k+1}$. In total, the cost of creating $Z_1 X_k W_{\geq k+1}$ from $Z_1 W_{\geq k+1}$ is $2k(k-1) - 4 + 2 + \sum_{j=2}^{k-1} 2 = 2k^2-4$ and the cost of creating $Z_1 Y_k W_{\geq k+1}$ from $Z_1 W_{\geq k+1}$ is $2k^2-5$.

This cost calculation does not change if the leading symbols are identities. In that scenario, we can use a modified version of the third subroutine, where we start with
\begin{equation}
    I_1 X_2 Z_3 Z_4 \cdots Z_k W_{\geq k+1}.
\end{equation}
Applying then the commutators with $X_j X_{j+1}$ and $Z_{j+1}$ for $ k = 2, \cdots, k-1$ results in
\begin{equation}
    I_1 I_2 \cdots I_{k-1} X_k W_{\geq k+1}.
\end{equation}

Note that this approach is not necessarily the most efficient one in how the commutators are applied at the first 2 indices, but gives a bound on the necessary number of commutators. The most expensive Pauli string under this algorithm consists of $X_1 X_2 \cdots X_n$ and costs
\begin{equation}
    2+\sum_{k=2}^n (2k^2-4) = \frac23 n^3 +O(n^2).
\end{equation}
See Appendix \ref{app:alg-computation} for a pseudocode writeup of the described approach.\end{proof}

\end{document}
